% Blueprint: rigorous error bound for Hardy's Z function via arbitrary-order
% Euler--Maclaurin for the Gamma and zeta components.

\chapter{Rigorous error bound for Hardy's $Z$ function}
\label{cha:hardy_z}

\section{The Euler--Maclaurin expansion of \texorpdfstring{$\log\Gamma$}{log Gamma}}
\label{sec:hz_loggamma}

Throughout this section $w\in\mathbb{C}$ with $\operatorname{Re}w>0$, and
$\operatorname{Log}\Gamma$ denotes the holomorphic branch of $\log\Gamma$ on the
right half-plane $\Omega:=\{\operatorname{Re}>0\}$: since $\Omega$ is simply connected
and $\Gamma$ is holomorphic and nowhere zero there, $\log\Gamma$ has a holomorphic
branch on $\Omega$, unique once normalised to be real on the positive real axis
($\operatorname{Log}\Gamma(x)=\log\Gamma(x)\in\mathbb{R}$ for $x>0$). Its derivative is
the digamma function $\psi=\Gamma'/\Gamma$, and its imaginary part is precisely the
continuous $\arg\Gamma$ (anchored at $\arg\Gamma(\tfrac14)=0$) used to define $\theta$
in \Cref{def:nf_theta}.

Recall from \Cref{thm:gs_digamma_stirling} that for every $M\in\mathbb{N}$,
\begin{equation}
\tag{$\psi$}
\label{eq:hz_digamma}
  \psi(w) \;=\; \operatorname{Log}w-\frac1{2w}
    -\sum_{k=1}^{M}\frac{B_{2k}}{2k}\,w^{-2k}
    \;-\; R^{\psi}_{2M+1}(w),
  \qquad R^{\psi}_{2M+1}(w):=\texttt{digammaRemainderIoi}\,(w,2M+1),
\end{equation}
with the explicit bound (\Cref{thm:gs_remainder_bound_explicit}, valid $\operatorname{Re}w>0$,
$2M+1\ge2$, i.e. $M\ge1$)
\begin{equation}
\tag{$\psi$-bd}
\label{eq:hz_digamma_bound}
  \bigl|R^{\psi}_{2M+1}(w)\bigr|\;\le\;\frac{(2M+1)!}{(2\pi)^{2M}}\,\lvert w\rvert^{-(2M+1)}.
\end{equation}

\begin{definition}[Stirling approximant of $\log\Gamma$]
\label{def:hz_stirling}
\lean{Complex.Stirling.logGammaStirling}
\leanok
\uses{def:bernoulli}
For $M\in\mathbb{N}$ and $\operatorname{Re}w>0$ set
\[
  S_M(w)\;:=\;\Bigl(w-\tfrac12\Bigr)\operatorname{Log}w-w+\tfrac12\log(2\pi)
    +\sum_{k=1}^{M}\frac{B_{2k}}{2k(2k-1)}\,w^{\,1-2k}.
\]
\end{definition}

\begin{lemma}[$S_M$ is a primitive of the digamma partial sum]
\label{lem:hz_stirling_deriv}
\lean{Complex.Stirling.hasDerivAt_logGammaStirling}
\leanok
\uses{def:hz_stirling}
For $\operatorname{Re}w>0$, $S_M$ is holomorphic and
\[
  S_M'(w)\;=\;\operatorname{Log}w-\frac1{2w}-\sum_{k=1}^{M}\frac{B_{2k}}{2k}\,w^{-2k},
\]
i.e. $S_M'$ is the truncated digamma series of \textup{($\psi$)}.
\end{lemma}

\begin{proof}
\leanok
Term by term. $\tfrac{d}{dw}\bigl[(w-\tfrac12)\operatorname{Log}w-w\bigr]
=\operatorname{Log}w+(w-\tfrac12)\tfrac1w-1=\operatorname{Log}w-\tfrac1{2w}$; the
constant $\tfrac12\log(2\pi)$ differentiates to $0$; and for $k\ge1$,
$\tfrac{d}{dw}\bigl[\tfrac{B_{2k}}{2k(2k-1)}w^{1-2k}\bigr]
=\tfrac{B_{2k}}{2k(2k-1)}(1-2k)w^{-2k}=-\tfrac{B_{2k}}{2k}w^{-2k}$. Summing gives
the claim; holomorphy on $\{\operatorname{Re}>0\}$ is clear ($\operatorname{Log}$ and
the powers are holomorphic there).
\end{proof}

\begin{definition}[Improper $\log\Gamma$ remainder]
\label{def:hz_remainder}
\lean{Complex.Stirling.logGammaRemainderIoi}
\leanok
\uses{def:gs_remainder_Ioi}
For $\operatorname{Re}w>0$ and $M\ge1$ define
\[
  \rho_M(w)\;:=\;\int_0^{\infty} R^{\psi}_{2M+1}(w+\tau)\,d\tau
    \;=\;\int_0^{\infty}\texttt{digammaRemainderIoi}\,(w+\tau,\,2M+1)\,d\tau .
\]
The integrand is continuous in $\tau$ and, by \textup{($\psi$-bd)},
$\bigl|R^{\psi}_{2M+1}(w+\tau)\bigr|\le (2M+1)!\,(2\pi)^{-2M}(\operatorname{Re}w+\tau)^{-(2M+1)}$
which is integrable on $[0,\infty)$ (as $2M+1\ge2$), so $\rho_M(w)$ is well defined.
\end{definition}

\begin{theorem}[Explicit arbitrary-order $\log\Gamma$ remainder bound]
\label{thm:hz_loggamma_bound}
\lean{Complex.Stirling.norm_logGammaRemainderIoi_le}
\leanok
\uses{def:hz_remainder, thm:gs_remainder_bound_explicit}
For every $M\ge1$ and every $w$ with $\operatorname{Re}w>0$,
\[
  \bigl|\rho_M(w)\bigr|\;\le\;\frac{(2M+1)!}{2M\,(2\pi)^{2M}}\,(\operatorname{Re}w)^{-2M}.
\]
In particular, for fixed $\operatorname{Im}w$, $\rho_M(w)\to0$ as $\operatorname{Re}w\to+\infty$.
\end{theorem}

\begin{proof}
\leanok
\uses{thm:gs_remainder_bound_explicit}
By the triangle inequality for the Bochner integral and
\textup{($\psi$-bd)}, using $\lvert w+\tau\rvert\ge\operatorname{Re}w+\tau>0$,
\[
  \bigl|\rho_M(w)\bigr|
  \le\int_0^\infty\bigl|R^{\psi}_{2M+1}(w+\tau)\bigr|\,d\tau
  \le\frac{(2M+1)!}{(2\pi)^{2M}}\int_0^\infty(\operatorname{Re}w+\tau)^{-(2M+1)}\,d\tau .
\]
With $\sigma:=\operatorname{Re}w>0$ and $2M+1\ge2$,
$\int_0^\infty(\sigma+\tau)^{-(2M+1)}d\tau=\sigma^{-2M}/(2M)$, giving the bound.
\end{proof}

The additive Stirling constant $\tfrac12\log(2\pi)$ in \Cref{def:hz_stirling} cannot be
recovered from the digamma expansion alone (differentiation kills constants). It is
pinned by the classical real-axis Stirling limit, our sole external analytic input in
this section (exactly analogous to the role of the vanishing digamma constant,
\Cref{lem:gs_anchor}, one derivative down).

\begin{lemma}[Real-axis Stirling anchor]
\label{lem:hz_anchor}
\lean{Real.Stirling.tendsto_log_Gamma_sub_stirlingRef_atTop}
\leanok
For real $x\to+\infty$,
\[
  \log\Gamma(x)\;-\;\Bigl[(x-\tfrac12)\log x-x+\tfrac12\log(2\pi)\Bigr]\;\longrightarrow\;0 .
\]
\end{lemma}

\begin{proof}
\leanok
This is Stirling's approximation for the Gamma function on the positive real axis. It
follows from mathlib's factorial Stirling equivalence
\texttt{Stirling.factorial\_isEquivalent\_stirling}
($n!\sim\sqrt{2\pi n}\,(n/e)^n$, i.e.
$\log n!-[(n+\tfrac12)\log n-n+\tfrac12\log(2\pi)]\to0$) together with the convexity
of $\log\circ\Gamma$ on $(0,\infty)$ (mathlib \texttt{convexOn\_log\_Gamma} in the
\texttt{Real} namespace). Convexity gives \emph{monotone difference quotients}: for
$x\in[n,n+1]$,
\[
  \log\Gamma(n)+ (x-n)\bigl(\log\Gamma(n)-\log\Gamma(n-1)\bigr)
  \;\le\;\log\Gamma(x)\;\le\;
  \log\Gamma(n)+(x-n)\bigl(\log\Gamma(n+1)-\log\Gamma(n)\bigr),
\]
i.e. $\log\Gamma(x)$ is trapped between the tangent-from-the-left and the chord, both of
which equal $\log\Gamma(n)+O(1)\cdot(x-n)$ with the $O(1)$ slopes
$\log\Gamma(n{\pm}1)-\log\Gamma(n)=\log n+o(1)$. Subtracting the smooth reference
$r(x):=(x-\tfrac12)\log x-x+\tfrac12\log(2\pi)$ (for which $r(x)-r(n)\to0$ uniformly on
$[n,n+1]$ since $r'(x)=\log x-\tfrac1{2x}$ varies by $o(1)$ across the unit interval) and
using $\log\Gamma(n)-r(n)\to0$ (the factorial equivalence, as $\log\Gamma(n)=\log(n-1)!$)
squeezes $\log\Gamma(x)-r(x)\to0$ for all real $x\to+\infty$.
\end{proof}

\begin{theorem}[Euler--Maclaurin/Stirling expansion of $\log\Gamma$]
\label{thm:hz_loggamma_em}
\lean{Complex.Stirling.logGamma_eq_stirling_add_remainder}
\leanok
\uses{def:hz_stirling, def:hz_remainder, lem:hz_stirling_deriv, thm:hz_loggamma_bound,
      lem:hz_anchor, thm:gs_digamma_stirling}
For every $M\ge1$ and every $w$ with $\operatorname{Re}w>0$,
\[
  \operatorname{Log}\Gamma(w)\;=\;S_M(w)\;+\;\rho_M(w).
\]
\end{theorem}

\begin{proof}
\leanok
\uses{lem:hz_stirling_deriv, thm:gs_digamma_stirling, thm:hz_loggamma_bound, lem:hz_anchor}
Put $g(w):=\operatorname{Log}\Gamma(w)-S_M(w)$, holomorphic on the right half-plane
$\Omega:=\{\operatorname{Re}>0\}$.

\emph{Shared derivative.} Digamma is the logarithmic derivative of $\Gamma$, so
$(\operatorname{Log}\Gamma)'=\psi$; with \Cref{lem:hz_stirling_deriv} and \textup{($\psi$)},
\[
  g'(w)=\psi(w)-S_M'(w)=-R^{\psi}_{2M+1}(w).
\]

\emph{Integrating the derivative along the horizontal ray.} We integrate $g'$, which uses
only the bound on $R^{\psi}$ itself. Since $t\mapsto g(w+t)$ has derivative $-R^{\psi}_{2M+1}(w+t)$ on $[0,n]$ (all
points have real part $>0$), and $R^{\psi}_{2M+1}(w+\cdot)$ is integrable on $(0,\infty)$ (its norm
is dominated by $C(\operatorname{Re}w+\tau)^{-(2M+1)}$ via \Cref{thm:hz_loggamma_bound}), the
fundamental theorem of calculus gives, for every $n\in\mathbb{N}$,
\[
  g(w)=\int_0^n R^{\psi}_{2M+1}(w+\tau)\,d\tau+g(w+n).
\]
As $n\to\infty$ the truncated integral converges to $\rho_M(w)=\int_0^\infty R^{\psi}_{2M+1}(w+\tau)\,d\tau$
(partial to improper integral). It remains to show $g(w+n)\to0$.

\emph{The shifted value vanishes.} Split
$g(w+n)=\bigl(g(w+n)-g(\operatorname{Re}w+n)\bigr)+g(\operatorname{Re}w+n)$. For the real point,
$g(\operatorname{Re}w+n)=\log\Gamma(\operatorname{Re}w+n)-S_M(\operatorname{Re}w+n)\to0$ by
\Cref{lem:hz_anchor} (the additive Bernoulli tail $\sum_{k=1}^{M}\tfrac{B_{2k}}{2k(2k-1)}x^{1-2k}\to0$).
For the difference, integrate $g'=-R^{\psi}\cdot i$ along the vertical segment from
$\operatorname{Re}w+n$ to $w+n$: by the fundamental theorem of calculus and
$\|R^{\psi}_{2M+1}(s)\|\le C\|s\|^{-(2M+1)}\le C(\operatorname{Re}w+n)^{-(2M+1)}$,
\[
  \bigl\|g(w+n)-g(\operatorname{Re}w+n)\bigr\|
  \le C\,(\operatorname{Re}w+n)^{-(2M+1)}\,\lvert\operatorname{Im}w\rvert\longrightarrow0.
\]
Hence $g(w+n)\to0$, and therefore $g(w)=\rho_M(w)$, which is the claim.
\end{proof}

\begin{remark}[Sharper decay]
\label{rem:hz_sharp}
Repeating the direct centered-kernel/P\'olya argument of \Cref{thm:zeta_error_bound}
for the kernel $\log(w+x)$ yields the
one-power-sharper bound $|\rho_M(w)|\le C_M\lvert w\rvert^{-(2M+1)}$. The
$(\operatorname{Re}w)^{-2M}$ form above is weaker by one factor of $|w|$; it follows
directly from the digamma bound, and suffices for
the algorithm, whose shift below makes $\operatorname{Re}w$ as large as desired.
\end{remark}

\section{The error bound for the Riemann--Siegel theta function}
\label{sec:hz_theta}

Write $w_T:=\tfrac14+i\tfrac T2$ (real part $\tfrac14>0$). The Riemann--Siegel theta
$\theta=\texttt{riemannSiegelTheta}$ of \Cref{def:nf_theta} is defined via the continuous
argument variation $\Delta_{\gamma_\Gamma}\arg\Gamma$ along $\gamma_\Gamma(u)=\tfrac14+i\tfrac u2$.
To use the $\log\Gamma$ expansion we first record the bridge to the holomorphic branch,
which is the content asserted in prose by \Cref{def:nf_theta}.

\begin{lemma}[Theta via the holomorphic $\log\Gamma$ branch]
\label{lem:hz_theta_log}
\lean{riemannSiegelTheta_eq_im_logGamma}
\leanok
\uses{def:nf_theta}
For every $T\ge0$,
\[
  \theta(T)\;=\;\operatorname{Im}\operatorname{Log}\Gamma(w_T)\;-\;\tfrac T2\log\pi ,
\]
where $\operatorname{Log}\Gamma$ is the holomorphic branch on $\Omega=\{\operatorname{Re}>0\}$
normalised real on the positive axis (\Cref{sec:hz_loggamma}); equivalently, the continuous
$\arg\Gamma$ along $\gamma_\Gamma$ anchored at $\arg\Gamma(\tfrac14)=0$ equals
$\operatorname{Im}\operatorname{Log}\Gamma(w_T)$.
\end{lemma}

\begin{proof}
\leanok
\uses{def:nf_theta}
The segment $\gamma_\Gamma([0,T])$ lies in $\Omega$, on which $\Gamma$ is holomorphic and
nowhere zero and $\operatorname{Log}\Gamma$ is a holomorphic primitive of $\psi$ with
$\operatorname{Im}\operatorname{Log}\Gamma$ real-analytic. By the defining property of the
continuous argument variation (\Cref{def:nf_argvar}), $\Delta_{\gamma_\Gamma}\arg\Gamma
=\operatorname{Im}\operatorname{Log}\Gamma(\gamma_\Gamma(T))-\operatorname{Im}\operatorname{Log}\Gamma(\gamma_\Gamma(0))$;
here $\gamma_\Gamma(T)=w_T$ and $\gamma_\Gamma(0)=\tfrac14$ with
$\operatorname{Im}\operatorname{Log}\Gamma(\tfrac14)=\arg\Gamma(\tfrac14)=0$ since
$\Gamma(\tfrac14)>0$. Subtracting $\tfrac T2\log\pi$ gives the claim (\Cref{def:nf_theta}).
\end{proof}

Its Euler--Maclaurin estimate truncated at order $M$ is
\[
  \widetilde\theta_M(T)\;:=\;\operatorname{Im}S_M(w_T)-\tfrac T2\log\pi .
\]

\begin{corollary}[Theta error, unshifted]
\label{cor:hz_theta_unshifted}
\lean{abs_riemannSiegelTheta_sub_stirling_le}
\leanok
\uses{cor:hz_theta_shifted}
For every $M\ge1$ and every $T\ge0$,
\[
  \bigl|\theta(T)-\widetilde\theta_M(T)\bigr|
  \;=\;\bigl|\operatorname{Im}\rho_M(w_T)\bigr|
  \;\le\;\bigl|\rho_M(w_T)\bigr|
  \;\le\;\frac{(2M+1)!}{2M\,(2\pi)^{2M}}\Bigl(\tfrac14\Bigr)^{-2M}
  \;=\;\frac{(2M+1)!}{2M}\Bigl(\tfrac2\pi\Bigr)^{2M}.
\]
\end{corollary}

\begin{proof}
\leanok
\uses{thm:hz_loggamma_em, thm:hz_loggamma_bound, lem:hz_theta_log}
Subtract, using \Cref{lem:hz_theta_log} for $\theta(T)$:
$\theta(T)-\widetilde\theta_M(T)=\operatorname{Im}\bigl(\operatorname{Log}\Gamma(w_T)-S_M(w_T)\bigr)
=\operatorname{Im}\rho_M(w_T)$ by \Cref{thm:hz_loggamma_em}. Then
$|\operatorname{Im}z|\le|z|$ and \Cref{thm:hz_loggamma_bound} at $w=w_T$ with
$\operatorname{Re}w_T=\tfrac14$ give the bound; finally
$(\tfrac14)^{-2M}=4^{2M}$ and $4^{2M}/(2\pi)^{2M}=(2/\pi)^{2M}$.
\end{proof}

To match the algorithm, which shifts the argument by an integer $N$ so that the
Stirling series converges rapidly, we use the exact functional recurrence of
$\Gamma$.

\begin{lemma}[Exact shift of $\log\Gamma$]
\label{lem:hz_shift}
\lean{Complex.logGamma_eq_sub_sum_log_add}
\leanok
For $\operatorname{Re}w>0$ and $N\in\mathbb{N}$,
\[
  \operatorname{Log}\Gamma(w)\;=\;\operatorname{Log}\Gamma(w+N)\;-\;\sum_{j=0}^{N-1}\operatorname{Log}(w+j).
\]
\end{lemma}

\begin{proof}
\leanok
Induction on $N$ from $\Gamma(z+1)=z\,\Gamma(z)$, i.e.
$\operatorname{Log}\Gamma(z+1)=\operatorname{Log}(z)+\operatorname{Log}\Gamma(z)$ on the
right half-plane (both sides holomorphic, agreeing up to $2\pi i\mathbb{Z}$, and equal
at large real $z$ by the anchoring of the branch). The base case $N=0$ is trivial;
the inductive step peels one factor.
\end{proof}

\begin{corollary}[Theta error, shifted: the algorithm's bound]
\label{cor:hz_theta_shifted}
\lean{abs_riemannSiegelTheta_sub_stirling_shift_le}
\leanok
\uses{lem:hz_shift, thm:hz_loggamma_em, thm:hz_loggamma_bound, lem:hz_theta_log}
Fix $N\in\mathbb{N}$, $M\ge1$ and $T\ge0$. Writing $w_T=\tfrac14+i\tfrac T2$, the shifted
order-$M$ Stirling estimate
\[
  \widetilde\theta_{M,N}(T)\;:=\;\operatorname{Im}\Bigl(S_M(w_T+N)-\sum_{j=0}^{N-1}\operatorname{Log}(w_T+j)\Bigr)-\tfrac T2\log\pi
\]
satisfies
\[
  \bigl|\theta(T)-\widetilde\theta_{M,N}(T)\bigr|
  \;=\;\bigl|\operatorname{Im}\rho_M(w_T+N)\bigr|
  \;\le\;\frac{(2M+1)!}{2M\,(2\pi)^{2M}}\,\bigl(\tfrac14+N\bigr)^{-2M}.
\]
\end{corollary}

\begin{proof}
\leanok
\uses{lem:hz_shift, thm:hz_loggamma_em, thm:hz_loggamma_bound, lem:hz_theta_log}
By \Cref{lem:hz_shift} the finite sum is exact, so
$\operatorname{Log}\Gamma(w_T)=\bigl[S_M(w_T+N)-\sum_{j=0}^{N-1}\operatorname{Log}(w_T+j)\bigr]+\rho_M(w_T+N)$
using \Cref{thm:hz_loggamma_em} at $w_T+N$. With \Cref{lem:hz_theta_log} for $\theta(T)$,
taking imaginary parts and subtracting
$\tfrac T2\log\pi$ shows $\theta(T)-\widetilde\theta_{M,N}(T)=\operatorname{Im}\rho_M(w_T+N)$.
Apply \Cref{thm:hz_loggamma_bound} at $w=w_T+N$, with
$\operatorname{Re}(w_T+N)=\tfrac14+N$.
\end{proof}

\section{The \texorpdfstring{$Z$}{Z}-function error bound}
\label{sec:hz_Z}

Hardy's $Z$ is the phase-rotated zeta value on the critical line; the algorithm
computes its real part (which equals $Z$ itself by the functional equation, though we
never need that reality for the error bound).

\begin{definition}[Hardy's $Z$ function]
\label{def:hz_Z}
\lean{Zeta.hardyZ}
\leanok
\uses{def:nf_theta}
For $t\in\mathbb{R}$ set
\[
  Z(t)\;:=\;\operatorname{Re}\bigl(e^{i\theta(t)}\,\zeta(\tfrac12+it)\bigr),
\]
with $\theta$ the Riemann--Siegel theta of \Cref{def:nf_theta}. This is the
real-valued quantity whose sign changes locate the zeros of $\zeta$ on the critical
line (\texttt{python/src/euler\_maclaurin/zeta\_zeros.py}).
\end{definition}

\begin{lemma}[Phase perturbation]
\label{lem:hz_phase}
\lean{norm_expImul_sub_expImul_le}
\leanok
For $\alpha,\beta\in\mathbb{R}$ and $z,w\in\mathbb{C}$,
\[
  \bigl\|e^{i\alpha}z-e^{i\beta}w\bigr\|\;\le\;\lvert\alpha-\beta\rvert\,\|z\|+\|z-w\|.
\]
\end{lemma}

\begin{proof}
\leanok
Write $e^{i\alpha}z-e^{i\beta}w=(e^{i\alpha}-e^{i\beta})z+e^{i\beta}(z-w)$.
Since $e^{i\alpha}-e^{i\beta}=e^{i\beta}(e^{i(\alpha-\beta)}-1)$ and $|e^{i\beta}|=1$,
\[
  \bigl\|e^{i\alpha}-e^{i\beta}\bigr\|=\bigl\|e^{i(\alpha-\beta)}-1\bigr\|\le\lvert\alpha-\beta\rvert,
\]
the last step being $\|e^{ix}-1\|=\bigl\|\int_0^x i\,e^{iu}\,du\bigr\|\le|x|$ (mathlib's
\texttt{Complex.norm\_exp\_ofReal\_mul\_I\_sub\_one\_le}). Combining
with $\|e^{i\beta}(z-w)\|=\|z-w\|$ and the triangle inequality gives the claim.
\end{proof}

\begin{theorem}[Rigorous $Z$-function error bound]
\label{thm:hz_Z_bound}
\lean{abs_hardyZ_sub_re_le}
\leanok
\uses{def:hz_Z, lem:hz_phase, cor:hz_theta_shifted, thm:zeta_error_bound}
Let $t\in\mathbb{R}$, and let $\widetilde\theta\in\mathbb{R}$, $\widetilde\zeta\in\mathbb{C}$
be any approximations of $\theta(t)$ and $\zeta(\tfrac12+it)$. Put
$\widetilde Z:=\operatorname{Re}\bigl(e^{i\widetilde\theta}\,\widetilde\zeta\bigr)$. Then
\[
  \bigl|Z(t)-\widetilde Z\bigr|
  \;\le\;\bigl\lvert\theta(t)-\widetilde\theta\bigr\rvert\;\bigl\|\zeta(\tfrac12+it)\bigr\|
    \;+\;\bigl\|\zeta(\tfrac12+it)-\widetilde\zeta\bigr\|.
\]
In particular, for $t\ge0$, taking $\widetilde\theta=\widetilde\theta_{M,N}(t)$ from
\Cref{cor:hz_theta_shifted} and $\widetilde\zeta$ the order-$M'$ Euler--Maclaurin zeta
value at $s=\tfrac12+it$ with cutoff $N'$ (so that
$\zeta(\tfrac12+it)-\widetilde\zeta=R_{M'}(s,N')$ is the zeta remainder of
\Cref{thm:zeta_em}), the two error factors are bounded explicitly by
\Cref{cor:hz_theta_shifted} and \Cref{thm:zeta_error_bound} respectively:
\[
  \bigl|Z(t)-\widetilde Z\bigr|
  \;\le\;\frac{(2M+1)!}{2M\,(2\pi)^{2M}}\bigl(\tfrac14+N\bigr)^{-2M}\,\bigl\|\zeta(\tfrac12+it)\bigr\|
    \;+\;\Bigl|\tfrac{s+2M'+1}{\operatorname{Re}s+2M'+1}\Bigr|\,\lvert t_{M'+1}\rvert,
\]
and $\|\zeta(\tfrac12+it)\|\le\|\widetilde\zeta\|+\|R_{M'}(s,N')\|$ makes the whole
bound computable from the approximants alone.
\end{theorem}

\begin{proof}
\leanok
\uses{def:hz_Z, lem:hz_phase, cor:hz_theta_shifted, thm:zeta_error_bound}
Since taking real parts is $1$-Lipschitz, $|\operatorname{Re}a-\operatorname{Re}b|\le\|a-b\|$,
\[
  \bigl|Z(t)-\widetilde Z\bigr|
  =\bigl|\operatorname{Re}\bigl(e^{i\theta(t)}\zeta(\tfrac12+it)\bigr)-\operatorname{Re}\bigl(e^{i\widetilde\theta}\widetilde\zeta\bigr)\bigr|
  \le\bigl\|e^{i\theta(t)}\zeta(\tfrac12+it)-e^{i\widetilde\theta}\widetilde\zeta\bigr\|.
\]
Apply \Cref{lem:hz_phase} with $\alpha=\theta(t)$, $\beta=\widetilde\theta$,
$z=\zeta(\tfrac12+it)$, $\zeta=\widetilde\zeta$ to obtain the first displayed bound.
Substituting the explicit theta bound (\Cref{cor:hz_theta_shifted}) for
$|\theta(t)-\widetilde\theta|$ and the explicit zeta remainder bound
(\Cref{thm:zeta_error_bound}) for $\|\zeta(\tfrac12+it)-\widetilde\zeta\|$ gives the
second; the final inequality is the triangle inequality
$\|\zeta(\tfrac12+it)\|\le\|\widetilde\zeta\|+\|\zeta(\tfrac12+it)-\widetilde\zeta\|$.
\end{proof}

\begin{remark}[Scope]
\label{rem:hz_scope}
\Cref{thm:hz_Z_bound} is the analytic error certificate for the value $Z(t)$;
interval/ball arithmetic on top of it turns the two explicit error factors into a
machine-checked bound and a certified sign of $Z(t)$. Reality of $Z$ (via the functional equation) and the
sign-change zero-localisation are downstream of this bound and are not needed for it.
\end{remark}
